\unnumberedchapter{ВВЕДЕНИЕ}

Онлайн-сервисы для шахмат давно перестали быть только виртуальной доской. Современная шахматная платформа должна объединять регистрацию пользователей, хранение профилей, рейтинги, быстрый обмен событиями, проверку корректности ходов, историю партий, анализ позиций, турнирные сценарии и удобный интерфейс для администратора. Из-за этого разработка такого приложения хорошо показывает практические навыки full-stack разработки: проектирование базы данных, создание API, работу с асинхронным сервером, организацию frontend-архитектуры и подготовку воспроизводимого запуска.

Актуальность дипломного проекта ChessView связана с тем, что шахматные онлайн-платформы требуют согласованной работы нескольких технических слоев. Пользователь ожидает мгновенный отклик интерфейса, но фактическое состояние партии должно подтверждаться сервером. Если клиент самостоятельно решал бы, законен ли ход и кто победил, система была бы уязвима к ошибкам и манипуляциям. Поэтому в ChessView используется server-authoritative подход: браузер отвечает за отображение и ввод действий, а backend проверяет правила игры, обновляет часы, сохраняет ход и рассылает новое состояние участникам.

Целью дипломного проекта является разработка веб-приложения ChessView, предназначенного для организации шахматных онлайн-партий, анализа позиций, ведения профилей, работы с задачами, турнирами, запланированными матчами и административным контролем пользователей.

Для достижения цели были поставлены следующие задачи: изучить проблематику шахматных онлайн-приложений; сравнить существующие решения; выбрать инструменты разработки; определить требования к серверу; спроектировать структуру программного обеспечения; разработать структуру базы данных; описать REST API и WebSocket API; реализовать клиентскую часть; реализовать административную панель; описать пользовательские и административные сценарии.

Объектом разработки является шахматное веб-приложение. Предметом разработки являются программные модули ChessView: backend на FastAPI, клиентская часть на React и TypeScript, административный frontend, база данных PostgreSQL, WebSocket-диспетчер, модуль шахматной логики, модуль турниров, модуль задач и административная подсистема.

Практическая значимость проекта заключается в создании воспроизводимой учебной платформы, которую можно запустить локально через Docker Compose или раздельно в режиме разработки. Результатом работы является программный продукт, пояснительная записка и Overleaf-проект, подготовленный с учетом требований к оформлению дипломного проекта.
